\begin{lemma}[Composition of replacement presentations]
Fix a finite loopless multigraph \(G\), a palette of \(k\) colors, a
distinguished edge \(e_0\), an input coloring, an exact ordered terminal
tree \(T\), a generating coloring \(\varphi\), an extended Tashkinov tree
\(\mathcal E\), a hierarchy base, a
\noderef{HierarchyCounterexample} \(C\), and its
\noderef{HierarchyMinimumCounterexampleContext}.  Let
\(P\) be a \noderef{HierarchyReplacementPresentation} relative to these
fixed original data.  Write \(\mu=P.\mathit{coloring}\) and
\(\mathcal E_\mu=P.\mathit{corresponding}.\mathit{replacementETT}\).

Let \(\tau\) be another coloring.  Suppose there is a
\noderef{CorrespondingExtendedTashkinovTree} from \(\mathcal E_\mu\) to
\(\tau\) on the same exact \(T\), a hierarchy base on its replacement ETT,
a hierarchy counterexample \(C_\tau\) on that base, and
\[
 \noderef{SameGoodHierarchy}
   (P.\mathit{counterexample}.\mathit{good},C_\tau.\mathit{good}).
\]
Then there is a replacement presentation \(P_\tau\), relative to the
original \(\mathcal E,C\), and the original minimum context, whose coloring
is exactly \(\tau\).
\end{lemma}
\begin{proof}
Put
\[
 \mathcal C_0=P.\mathit{corresponding},\qquad
 \mathcal C_1=\mathit{second},\qquad
 C_\mu=P.\mathit{counterexample}.
\]
Thus \(\mathcal C_0\) corresponds from the original ETT \(\mathcal E\) to
the intermediate ETT \(\mathcal E_\mu\), whereas \(\mathcal C_1\)
corresponds from \(\mathcal E_\mu\) to its replacement ETT
\(\mathcal E_\tau\).  Apply
\noderef{CorrespondingExtendedTashkinovTreeTrans} to
\(\mathcal C_0,\mathcal C_1\).  It supplies a composite certificate
\[
 \mathcal C_{10}:
 \noderef{CorrespondingExtendedTashkinovTree}(\mathcal E,\tau,T)       \tag{1}
\]
and three comparison witnesses: equality of its replacement input with
that of \(\mathcal C_1\), heterogeneous equality of the two dependent
replacement series, and heterogeneous equality of the two replacement
ETTs.

Use equality elimination on the replacement-input equality first.  Then
use dependent heterogeneous-equality elimination on the replacement-series
equality and on the replacement-ETT equality, in that order.  Transport
along these eliminations the supplied second base and \(C_\tau\).  By
\noderef{HierarchyCounterexample}, the latter transport includes
\(C_\tau\)'s positive-rung certificate, hierarchy, good-hierarchy witness,
maximum-property certificate, and non-elementarity certificate; transport
the stated same-good witness along with them.
After these transports, all these objects are indexed by the literal
replacement input and replacement ETT of (1).  In particular, no object
whose type merely mentions a heterogeneously equal ETT is reused without
transport.  We retain the notation \(\mathcal B_\tau,C_\tau\) for the
transported base and counterexample.

We next compose the two same-good witnesses.  The hierarchy and
good-hierarchy fields used here are precisely the fields packaged by
\noderef{HierarchyCounterexample} in \(C_\mu\) and \(C_\tau\).  Unfold
\noderef{SameGoodHierarchy} for \(P.\mathit{sameGood}\) and for the
transported second witness.  Their first clauses are two successive
\noderef{SameHierarchyStages} relations.  Transitivity of equality gives
equality of the first and third base cores and divider counts;
transitivity of heterogeneous equality gives equality of their dependent
stage families and divider-position families.  Compose the two
equivalences asserting that the terminal extension kind is PE.  If this
common assertion holds, apply the first closure equality and then the
second closure equality.  These are exactly all clauses of
\noderef{SameHierarchyStages} between \(C.\mathit{hierarchy}\) and
\(C_\tau.\mathit{hierarchy}\).

Let \(h_0:q=q_\mu\) and \(h_1:q_\mu=q_\tau\) be the two divider-count
equalities, and use \(h_0h_1:q=q_\tau\) for the composite.  Given
\(j<q+1\), cast it first along \(h_0\), then along \(h_1\).  Functoriality
of dependent casts identifies this iterated cast with the direct cast
along \(h_0h_1\).  The first same-good witness identifies the two
\noderef{HierarchyStoredColorsAt} sets at the first two indices, and the
second identifies the sets at the last two; equality transitivity gives
the required first-to-third stored-set equality.  If
\(\eta\) belongs to the first stored set, transport that membership across
the first equality, apply the second Gamma-pair equality to \(\eta\), and
compose it with the first Gamma-pair equality.  After the same cast
identification this proves equality of the first and third Gamma pairs.
Consequently we have constructed
\[
 s_{10}:\noderef{SameGoodHierarchy}
       (C.\mathit{good},C_\tau.\mathit{good}).                 \tag{2}
\]

It remains to build the minimum context on the data in (1).  We verify its
six fields separately, always starting with
\(P.\mathit{minimumContext}\).

First let \(n_\mu=\mathcal E_\mu.\mathit{rounds}\).  The terminal
properness field of \(\mathcal C_1\), at
\(\operatorname{Fin.last}(n_\mu)\), says that the terminal replacement
series coloring is proper.  Its final-generating-coloring equality
identifies that coloring with \(\tau\), so \(\tau\) is proper.  The
terminal exact-one-uncolored field and the same equality show that exactly
\(e_0\) is uncolored under \(\tau\).  These give
\(\mathit{proper}\) and \(\mathit{oneEdgeUncolored}\).

For \(\mathit{lowerRung}\), take any positive maximum ETT \(\mathcal L\)
whose rung number is strictly below that of \(\mathcal E_\tau\).  The
rung-agreement field of \(\mathcal C_1\), retained by the transports
above, is
\[
 \mathcal E_\tau.\mathit{rounds}
   =\mathcal E_\mu.\mathit{rounds}.                            \tag{3}
\]
Rewrite the strict inequality by (3) and apply the complete lower-rung
field of \(P.\mathit{minimumContext}\).  This transports the entire
six-assertion package, with its arbitrary input, tree, and generating
coloring, not just its elementarity clause.

For \(\mathit{fewerDividers}\), let \(\mathcal L\) be a maximum ETT with
the same rung as \(\mathcal E_\tau\), and let \(K\) be a good hierarchy
on it with fewer dividers than
\(C_\tau.\mathit{hierarchy}\).  Rewrite the rung equality by (3).  The
second supplied same-good witness gives
\[
 C_\mu.\mathit{hierarchy}.\mathit{dividerCount}
 =C_\tau.\mathit{hierarchy}.\mathit{dividerCount}.             \tag{4}
\]
Rewrite the strict divider inequality by (4).  The
\(\mathit{fewerDividers}\) field of \(P.\mathit{minimumContext}\) now
applies and proves the required elementarity.

For lexicographic minimality, we first prove the literal equality
\[
 \noderef{HierarchyCounterexampleMeasure}(C_\mu)
   =\noderef{HierarchyCounterexampleMeasure}(C_\tau).          \tag{5}
\]
The \noderef{SameHierarchyStages} part of the second same-good witness
gives (4), equality of the base cores, and transported equality of every
ordinary hierarchy stage.  Hence the last
\noderef{HierarchyStarredStage} is the same exact ordered tree on both
sides: at divider count zero it is the equal base core, and at positive
divider count it is the corresponding equal ordinary stage.  Both
hierarchies have the same exact terminal tree \(T\).  Expanding
\noderef{HierarchyPathNumber}, their tail lengths, finite candidate sets,
and candidate predicates are therefore identical, including the
empty-candidate fallback.  This proves equality of the first coordinates
in (5).  For the second coordinates, (4) and equality of the ordinary
stage families identify the stages at index equal to the divider count;
subtracting their equal lengths from the common number
\(T.\mathit{length}\) gives equality.  Equality of both coordinates
proves (5).

Now take any hierarchy counterexample \(C'\) with the same rung as
\(\mathcal E_\tau\) and the same divider count as \(C_\tau\).  Equations
(3) and (4) turn these into the corresponding equalities with
\(\mathcal E_\mu,C_\mu\).  If
\(\noderef{HierarchyCounterexampleLess}\) made the measure of \(C'\)
strictly smaller than the measure of \(C_\tau\), rewriting by (5) would
make it strictly smaller than the measure of \(C_\mu\), contradicting
the \(\mathit{lexicographicMinimality}\) field of
\(P.\mathit{minimumContext}\).  This proves the new minimality field.

Finally consider \(\mathit{largeMissingSet}\).  By
\noderef{HierarchyCounterexample}, the positive-rung field of \(C_\tau\)
says that \(\mathcal E_\tau\) has positive rung.  The rung-agreement field
of \(\mathcal C_{10}\) therefore makes
\(n=\mathcal E.\mathit{rounds}\) positive and permits the terminal positive
index \(i<n\) whose successor is \(\operatorname{Fin.last}(n)\).  Put
\[
 L=\mathcal E.\mathit{series}.\mathit{data}.\mathit{tree}
       (\operatorname{Fin.last}(n)).
\]
Apply at \(i\) the replacement-stability fields of \(\mathcal C_0\) and
\(\mathcal C_{10}\).  We now make explicit the internal identifications
needed to read these two stability relations at the terminal rung.  For
each of the two certificates, the
\(\mathit{finalGeneratingColoring}\) field of
\noderef{CorrespondingExtendedTashkinovTree} identifies its replacement
series terminal coloring with \(\mu\), respectively with \(\tau\).  Its
\(\mathit{seriesPresentation}\) field identifies the replacement ETT's
internal series with that displayed replacement series, and its
\(\mathit{ladderTreesAgree}\) field at \(i\) identifies the actual terminal
ladder tree of \(\mathcal E_\mu\), respectively of \(\mathcal E_\tau\),
with the exact ordered tree \(L\).  Finally, the
\(\mathit{generatingColoring}\) field of
\noderef{ExtendedTashkinovTree} identifies the original series terminal
coloring with \(\varphi\).

Thus, after these four literal identifications, the vertexwise clause of
\noderef{StableColoring} for \(\mathcal C_0\) gives
\(\overline\mu(v)=\overline\varphi(v)\) at every vertex \(v\) of \(L\),
and the corresponding clause for \(\mathcal C_{10}\) gives
\(\overline\tau(v)=\overline\varphi(v)\) there.  Taking the unions over the
same displayed vertex set, exactly as in \noderef{TreeMissingColors}, gives
\[
 \overline\mu(T_{n,\mu})=\overline\varphi(L)
   =\overline\tau(T_{n,\tau}),
\]
where \(T_{n,\mu}\) and \(T_{n,\tau}\) are the actual terminal ladder
trees in the internal series of \(\mathcal E_\mu\) and
\(\mathcal E_\tau\).  Hence these sets have equal cardinalities.  The rung
agreements of \(\mathcal C_0\) and \(\mathcal C_{10}\) also identify both
rung numbers with \(n\), so the two bounds \(2n+11\) are identical.
Therefore the old \(\mathit{largeMissingSet}\) inequality in
\(P.\mathit{minimumContext}\) is exactly the required inequality under
\(\tau\).  Notice that the input, series, and ETT equalities returned by
\noderef{CorrespondingExtendedTashkinovTreeTrans} were already used above
to transport the dependent base, counterexample, and same-good data; those
outer equalities do not replace the internal
\(\mathit{finalGeneratingColoring}\), \(\mathit{seriesPresentation}\), and
\(\mathit{ladderTreesAgree}\) equalities used in this paragraph.

The six fields just constructed form a
\noderef{HierarchyMinimumCounterexampleContext} \(\mathit{ctx}_\tau\)
on the literal data in (1).  By
\noderef{HierarchyCounterexample}, \(C_\tau\) supplies on the transported
base exactly its positive-rung, hierarchy, good-hierarchy, maximum-property,
and non-elementarity fields.  Use \(C_\tau.\mathit{maximum}\) for the
separate maximum field of the replacement presentation; use \(C_\tau\)
itself, including its positive-rung, hierarchy, good-hierarchy, and
non-elementarity fields, for the counterexample field; and use (2) for the
original-relative same-good field.  Then
\[
 P_\tau=(\tau,\mathcal C_{10},C_\tau.\mathit{maximum},
          \mathcal B_\tau,C_\tau,s_{10},\mathit{ctx}_\tau)
\]
fills every field of \noderef{HierarchyReplacementPresentation}.  Its
coloring field is definitionally \(\tau\), which proves the conclusion.
\end{proof}
